\documentclass[a4paper]{article} 

\addtolength{\hoffset}{-2.25cm}
\addtolength{\textwidth}{4.5cm}
\addtolength{\voffset}{-3.25cm}
\addtolength{\textheight}{5cm}
\setlength{\parskip}{0pt}
\setlength{\parindent}{0in}

%----------------------------------------------------------------------------------------
%	PACKAGES AND OTHER DOCUMENT CONFIGURATIONS
%----------------------------------------------------------------------------------------
\usepackage{caption}
\usepackage{charter} % Use the Charter font
\usepackage[utf8]{inputenc} % Use UTF-8 encoding
\usepackage{microtype} % Slightly tweak font spacing for aesthetics
\usepackage[italian]{babel} % Language hyphenation and typographical rules
\usepackage{amsthm, amsmath, amssymb} % Mathematical typesetting
\usepackage{float} % Improved interface for floating objects
\usepackage[final, colorlinks = true, 
linkcolor = black, 
citecolor = black]{hyperref} % For hyperlinks in the PDF
\usepackage{graphicx, multicol} % Enhanced support for graphics
\usepackage{xcolor} % Driver-independent color extensions
\usepackage{rotating} % Rotation tools
\usepackage{censor} % Facilities for controlling restricted text
\usepackage{matlab-prettifier} % Environment for specifying algorithms in a natural way
\usepackage{booktabs} % Enhances quality of tables
\usepackage[yyyymmdd]{datetime} % Uses YEAR-MONTH-DAY format for dates
\usepackage[export]{adjustbox}
\renewcommand{\dateseparator}{-} % Sets dateseparator to '-'
\usepackage{fancyhdr} % Headers and footers
\pagestyle{fancy} % All pages have headers and footers
\fancyhead{}\renewcommand{\headrulewidth}{0pt} % Blank out the default header
\fancyfoot[L]{} % Custom footer text
\fancyfoot[C]{} % Custom footer text
\fancyfoot[R]{\thepage} % Custom footer text
\usepackage[nottoc,numbib]{tocbibind}

\begin{document}

%-------------------------------
%	TITLE SECTION
%-------------------------------

\fancyhead[C]{}
\hrule \medskip % Upper rule
\begin{minipage}{0.295\textwidth} 
\raggedright
\footnotesize
Filippo Falqui Cao \hfill\\   
699396 \hfill\\
f.falquicao@studenti.unipi.it
\end{minipage}
\begin{minipage}{0.4\textwidth} 
\centering 
\large 
Progetto Algoritmi e Strutture Dati\\ 
\normalsize 
Operazioni di ricerca minmax su grafi\\ 
\end{minipage}
\begin{minipage}{0.295\textwidth} 
\raggedleft
\today\hfill\\
\end{minipage}
\medskip\hrule 
\bigskip

%-------------------------------
%	CONTENTS
%-------------------------------

\tableofcontents

\bigskip\bigskip

\section{Linee Generali}
\subsection{Definizione del problema}
Il progetto ha l'obiettivo di analizzare la rete di trasmissione di pacchetti su internet tra Autonomous Systems. I dati derivano da misurazioni reali, come pubblicato dal dataset CAIDA AS Relationships. Da questo database occorre quindi costruire un grafo pesato, per poi saper rispondere a domande su coppie di vertici del grafo. 

Dai dati forniti da CAIDA, si possono estrarre una serie di righe, ciascuna che descrive un percorso con formato 

\medskip
{
\centering 
\texttt{routeviews/routeviews|5 1239|6113|8063|3464 199.88.21.0/24 i 144.228.240.93}\\
}
\medskip
Questa riga produce il percorso del pacchetto di informazioni $(1239,6113,8063,34641)$. Questa riga ha l'effetto sul grafo non diretto pesato di aumentare di $1$ il peso dell'arco $(1239,6113)$, di $(6113, 8063)$, e di $(8063, 34641)$. Una volta processate tutte le righe, ci limitiamo a considerare la componente connessa del grafo di dimensione massima. 

Dati due nodi $u$ e $v$ della componente connessa, consideriamo l'insieme (non vuoto) di tutti i percorsi $P$ che connettono $u$ e $v$. Ricordiamo che un percorso è una lista ordinata di nodi, ciascuno connesso da un arco con il successivo. Dunque $P_1 = u$ e $P_k = v$. Associamo a ogni percorso $P$ la funzione $\text{costo}(P)$:

$$
    \text{costo}(P) =  \max_{0\le i< k} w(P_i, P_{i+1})
$$

Dove $w$ è il peso dell'arco, mentre $k$ è la lunghezza del percorso.
L'obiettivo è di riuscire a trovare il minimo del costo di $P$, al variare di $P$ tra i due vertici $u$ e $v$ dati in input. 

\subsection{Piani preliminari}
Com'è chiaro dalla descrizione del problema, esso può essere comodamente suddiviso in due parti essenzialmente indipendenti: l'elaborazione del file CAIDA in un grafo, e l'elaborazione di un algoritmo per rispondere alle domande sul grafo. Ho preso quindi la decisione di creare due file distinti per ciascuna parte, usando il linguaggio \texttt{C++} per entrambi. Il primo file avrà quindi il compito di produrre una descrizione del grafo semplice e veloce da leggere per il secondo file. Siccome la maggior parte degli algoritmi sui grafi pesati può essere resa più veloce se i nodi sono interi consecutivi e gli archi vengono processati in ordine, ho deciso di assegnare l'incarico di riordinare gli archi e fornire un elenco preliminare di tutti i nodi al primo file, di modo da non dovermene più preoccupare nella seconda parte (quella concettualmente più difficile). Tutti gli step, per entrambi i file, vengono testati su \texttt{19980101.all-paths}, in quanto è il file più leggero disponibile. Ciononostante, dover leggere $43\cdot10^6$ caratteri di testo ogni volta non era un processo velocissimo, almeno sul mio computer; mi è quindi capitato di scrivere l'input direttamente da terminale a mano in fase di debugging. Per quanto riguarda il processo di documentazione, oltre ovviamente ai commenti all'interno del codice e alle descrizioni dei commit su Github, ho mantenuto un file \texttt{Appunti.txt} che conteneva le mie idee su come proseguire nello sviluppo. Questo file è stato invece prodotto a lavoro terminato.

\section{Setup}

\section{Prima soluzione}

\section{Seconda soluzione efficiente}
\subsection{Idea dell'algoritmo}
\subsection{Descrizione e correttezza}
\subsection{Implementazione}

\section{Punto facoltativo}


\bibliographystyle{plain} 
\bibliography{refs} 

\bigskip
\bigskip
\bigskip
\bigskip
\bigskip
\bigskip



{\bf Uso di GenAI:} nessuno.


\end{document}
